\documentclass[12pt,a4paper]{article}

% ── Packages ──────────────────────────────────────────────────────────────────
\usepackage[T1]{fontenc}
\usepackage[utf8]{inputenc}
\usepackage[margin=2.5cm]{geometry}
\usepackage{amsmath,amssymb,amsthm}
\usepackage{mathtools}
\usepackage{bm}
\usepackage{booktabs}
\usepackage{array}
\usepackage{xcolor}
\usepackage{graphicx}
\usepackage[ruled,vlined,linesnumbered]{algorithm2e}
\usepackage{microtype}
\usepackage{lmodern}
\usepackage[hidelinks,
            colorlinks=true,
            linkcolor=blue!60!black,
            citecolor=green!50!black,
            urlcolor=blue!70!black]{hyperref}
\usepackage{natbib}
\usepackage{setspace}

% ── Theorem environments ──────────────────────────────────────────────────────
\newtheorem{proposition}{Proposition}
\newtheorem{remark}{Remark}
\newtheorem{definition}{Definition}

% ── Notation macros ───────────────────────────────────────────────────────────
\newcommand{\obs}{\mathcal{T}}
\newcommand{\aug}{\mathcal{T}^+}
\newcommand{\brts}{\mathbf{t}}
\newcommand{\pars}{\bm{\theta}}
\newcommand{\lb}{\bm{\ell}}
\newcommand{\ub}{\mathbf{u}}
\newcommand{\fhat}{\widehat{f}}
\newcommand{\Lhat}{\widehat{L}}
\newcommand{\E}{\mathbb{E}}
\newcommand{\R}{\mathbb{R}}
\newcommand{\logf}{\ell^{(f)}}
\newcommand{\logq}{\ell^{(q)}}
\newcommand{\lw}{w}
\newcommand{\ESS}{\mathrm{ESS}}
\newcommand{\AIC}{\mathrm{AIC}}
\newcommand{\KL}{\mathrm{KL}}
\DeclareMathOperator*{\argmax}{arg\,max}
\DeclareMathOperator{\Var}{Var}
\newcommand{\dd}{\,\mathrm{d}}

% ── Title ─────────────────────────────────────────────────────────────────────
\title{%
  \textbf{\textsf{emphasis}}: Evolutionary Modeling for Phylogenetic
  Inference\\[0.4em]
  \large Technical Reference Manual}
\author{Francisco Richter\thanks{University of Southern Switzerland.
  Email: \texttt{richtf@usi.ch}}
  \and Thijs Janzen
  \and Hanno Hildenbrandt}
\date{\today}

% ══════════════════════════════════════════════════════════════════════════════
\begin{document}
\maketitle
\tableofcontents
\bigskip
\onehalfspacing

% ══════════════════════════════════════════════════════════════════════════════
\begin{abstract}
\noindent
\textsf{emphasis} is an \textsf{R} package for studying species diversification
and its ecological drivers. It implements a general birth--death model in which
per-lineage speciation and extinction rates depend on arbitrary state covariates:
current species richness, total phylogenetic diversity, and lineage-specific
evolutionary time. Because the marginal likelihood over the unobserved extinct
lineages is analytically intractable, the package uses \emph{importance sampling
on augmented trees}: a proposal distribution augments the observed extant tree
with stochastically drawn extinct lineages, and the resulting Monte Carlo
estimator of the marginal log-likelihood serves as the objective for two
complementary optimisers -- a Monte Carlo Cross-Entropy Method with simulated
annealing (CEM-SA) for global search and a Monte Carlo EM (MCEM) algorithm
for local refinement. Model selection uses the Akaike Information Criterion
computed from the IS log-likelihood at the optimum. All heavy computations are
parallelised at the particle level using shared-memory multithreading via
\textsf{RcppParallel}.
This document provides a self-contained mathematical derivation of every
component.
\end{abstract}

% ══════════════════════════════════════════════════════════════════════════════
\section{Notation}\label{sec:notation}

Table~\ref{tab:notation} summarises the principal symbols used throughout.

\begin{table}[ht]
\centering
\caption{Principal notation.}\label{tab:notation}
\renewcommand{\arraystretch}{1.25}
\begin{tabular}{lp{10cm}}
\toprule
Symbol & Meaning \\
\midrule
$\obs$ & Observed extant tree (phylo object or branching-time vector) \\
$z$ & Unobserved extinct lineages \\
$\aug = (\obs, z)$ & Augmented tree: $\obs$ together with extinct lineages $z$ \\
$\brts = (t_1 > \cdots > t_n > 0)$ & Ordered branching times of the extant tree \\
$\pars \in \R^8$ & Full model parameter vector \\
$\lambda(s,t;\pars)$ & Instantaneous speciation rate of lineage $s$ at time $t$ \\
$\mu(s,t;\pars)$ & Instantaneous extinction rate of lineage $s$ at time $t$ \\
$N(t)$ & Number of extant lineages at time $t$ \\
$P(t)$ & Total pendant branch-length at time $t$ \\
$E(s,t)$ & Evolutionary pendant time of lineage $s$ \\
$p(\obs, z \mid \pars)$ & Complete-data likelihood \\
$p(\obs \mid \pars)$ & Marginal (observed-data) likelihood \\
$q(z \mid \obs, \pars)$ & Proposal density for extinct lineages \\
$\logf_i$ & $\log p(\obs, z_i \mid \pars)$ \\
$\logq_i$ & $\log q(z_i \mid \obs, \pars)$ \\
$\lw_i$ & IS log-weight: $\lw_i = \logf_i - \logq_i$ \\
$S$ & IS sample size (augmented trees per particle) \\
$M$ & CEM population size (particles per iteration) \\
$\fhat(\pars)$ & IS estimate of $\log p(\obs \mid \pars)$ \\
$\ESS$ & Effective sample size \\
$\rho$ & Elite fraction retained per CEM iteration \\
$T$ & Maximum number of CEM iterations \\
$\bm{\sigma}^{(k)}$ & Per-parameter perturbation SD at CEM iteration $k$ \\
$\alpha$ & Annealing decrement per iteration \\
\bottomrule
\end{tabular}
\end{table}

% ══════════════════════════════════════════════════════════════════════════════
\section{The Diversification Model}\label{sec:model}

\subsection{General rate specification}

The \textsf{emphasis} model belongs to the family of state-dependent
birth--death processes \citep{etienne2012prolonged,janzen2015diversity}.
Every lineage $s$ alive at time $t$ (measured forward from the crown) is
subject to a per-lineage speciation rate and extinction rate that depend on
the current macroecological state of the clade:
\begin{align}\label{eq:rates}
  \lambda(s,t;\pars) &= f\!\left(\beta_0 + \beta_N N(t)
                         + \beta_P P(t) + \beta_E E(s,t)\right), \nonumber\\
  \mu(s,t;\pars)     &= f\!\left(\gamma_0 + \gamma_N N(t)
                         + \gamma_P P(t) + \gamma_E E(s,t)\right),
\end{align}
where the full parameter vector is
$\pars = (\beta_0, \beta_N, \beta_P, \beta_E,
          \gamma_0, \gamma_N, \gamma_P, \gamma_E)^\top \in \R^8$.

\paragraph{State covariates.}
\begin{description}
  \item[$N(t)$] \emph{Species richness.} The number of lineages currently
    alive. Negative $\beta_N$ produces diversity-dependent suppression of
    speciation (ecological carrying capacity), while positive $\beta_N$
    encodes positive diversity feedbacks.
  \item[$P(t)$] \emph{Phylogenetic diversity (PD).} The sum of all pendant
    branch-lengths:
    $P(t) = \sum_{s:\,\text{alive at }t} E(s,t)$.
    PD measures the ecological disparity available to the clade.
  \item[$E(s,t)$] \emph{Evolutionary pendant time.} The time elapsed since
    lineage $s$ last diverged. Long $E$ may indicate accumulated reproductive
    isolation (higher $\lambda$) or competitive displacement (higher $\mu$).
\end{description}

\paragraph{Link function.}
The link function $f:\R\to[0,\infty)$ ensures non-negative rates:
\begin{equation}\label{eq:link}
  f(\eta) =
  \begin{dcases}
    \max(0, \eta) & \text{linear link (default)}, \\
    \exp(\eta)    & \text{exponential link}.
  \end{dcases}
\end{equation}
The linear link allows rates to reach exactly zero and is piecewise-linear
in $\pars$; the exponential link guarantees strictly positive rates. The
exponential link is not supported for the EP covariate because no closed-form
integral exists for $\int \exp(\gamma_0 + \gamma_E(t-a_s))\dd t$ combined
with other covariates.

\subsection{Model shortcuts}\label{sec:models}

Setting subsets of $\{\beta_N, \beta_P, \beta_E, \gamma_N, \gamma_P,
\gamma_E\}$ to zero yields the four canonical models (Table~\ref{tab:models}).

\begin{table}[ht]
\centering
\caption{Model shortcuts and their free parameters.}\label{tab:models}
\renewcommand{\arraystretch}{1.2}
\begin{tabular}{lllp{7cm}}
\toprule
String & Formula & Active covariates & Free parameters \\
\midrule
\texttt{"cr"} & $\sim 1$ & none
  & $(\beta_0,\; \gamma_0)$ \\
\texttt{"dd"} & $\sim N$ & $N(t)$
  & $(\beta_0,\; \beta_N,\; \gamma_0,\; \gamma_N)$ \\
\texttt{"pd"} & $\sim \mathrm{PD}$ & $P(t)$
  & $(\beta_0,\; \beta_P,\; \gamma_0,\; \gamma_P)$ \\
\texttt{"ep"} & $\sim \mathrm{EP}$ & $E(s,t)$
  & $(\beta_0,\; \beta_E,\; \gamma_0,\; \gamma_E)$ \\
--- & $\sim N + \mathrm{PD}$ & $N(t),P(t)$
  & $(\beta_0,\beta_N,\beta_P,\gamma_0,\gamma_N,\gamma_P)$ \\
--- & $\sim N + \mathrm{PD} + \mathrm{EP}$ & all
  & all 8 parameters \\
\bottomrule
\end{tabular}
\end{table}

The number of free parameters is always $p = 2 + 2k$, where $k\in\{0,1,2,3\}$
is the number of active covariates.

% ══════════════════════════════════════════════════════════════════════════════
\section{Stochastic Simulation of Phylogenies}\label{sec:sim}

\subsection{Forward simulation}

Forward simulation of the state-dependent birth--death process proceeds via
Gillespie's direct method \citep{gillespie1977exact}. At each step the
algorithm:
\begin{enumerate}
  \item Computes the total event rate
    $R(t) = \sum_{s:\,\text{alive}} [\lambda(s,t;\pars) + \mu(s,t;\pars)]$.
  \item Draws the waiting time to the next event,
    $\Delta t \sim \mathrm{Exp}(R(t))$.
  \item Selects which lineage and event type (speciation or extinction)
    with probability proportional to individual rates.
  \item Updates the state $(N, P, \{E_s\})$ and the lineage set.
\end{enumerate}
Because $N(t)$, $P(t)$, and each $E(s,t)$ change only at events, $R(t)$
is piecewise-constant between events and the Gillespie step is exact.
For the EP covariate, where $E(s,t)$ grows continuously, the waiting time
no longer has a closed-form distribution and is drawn instead via
\emph{thinning} (Lewis--Shedler algorithm
\citep{lewis1979simulation}): candidate times are drawn from a
homogeneous Poisson process at rate $\lambda_{\max} \geq \sup_t R(t)$, and
each candidate is accepted with probability $R(t^*)/\lambda_{\max}$.

\subsection{Conditional simulation (augmentation)}\label{sec:proposal}

Given the observed extant tree $\obs$, the package draws the unobserved
extinct lineages $z$ from the proposal distribution $q(z \mid \obs, \pars)$,
which conditions on $\obs$ being the extant subtree. The algorithm operates
interval by interval between consecutive branching times
$(t_1 > t_2 > \cdots > t_n)$ of $\obs$:
\begin{enumerate}
  \item Within each interval $[t_{i+1}, t_i]$, candidate extinction events
    are generated by thinning a Poisson process at rate $\lambda_{\max}$.
  \item Each candidate time $t^*$ is accepted as an extinction event with
    probability $\lambda(t^*;\pars) / \lambda_{\max}$.
  \item An accepted extinction at $t^*$ spawns a new lineage, which is
    then forward-simulated from $t^*$ until it goes extinct before $t = 0$
    (the present); if it does not go extinct it is discarded and redrawn.
  \item The procedure recurses on the extinct branches.
\end{enumerate}
The log-proposal probability
\begin{equation}\label{eq:logq}
  \logq_i = \log q(z_i \mid \obs, \pars)
\end{equation}
is accumulated analytically during each simulation pass and returned
alongside $z_i$ and $\logf_i$, so that importance-sampling reweighting
can be applied immediately without re-running any simulation.

\subsection{Batch simulation and survival probability}

Passing a parameter matrix to \texttt{simulate\_tree()} triggers batch mode:
each row $\pars_r$ is simulated independently, with failed runs
(clade extinction or size overflow) retained. The fraction of simulations
that produce a surviving clade estimates the probability of non-extinction:
\begin{equation}
  \hat{\pi}(\pars) = \frac{1}{R}\sum_{r=1}^R
  \mathbf{1}[\text{simulation }r\text{ survived}].
\end{equation}
This quantity can be used to correct the likelihood for conditioning on
non-extinction. For diversity-dependent models, \texttt{train\_GAM()} fits a
binomial generalised additive model to the outcomes of a batch run, producing
a smooth estimate of $\pi(\pars)$ over the parameter space that can be
evaluated at any $\pars$ without further simulation.

\subsection{Parallel simulation}

The C++ simulation kernel is parallelised using \textsf{RcppParallel}, which
provides a portable thread pool backed by Intel TBB on supported platforms and
a fallback \texttt{std::thread} implementation elsewhere. Each particle (i.e.\
each augmented tree draw) is independent: particle $i$ runs its own
thinning loop, accumulates its own $(\logf_i, \logq_i)$ pair, and writes to
a pre-allocated per-particle slot with no shared mutable state. The
thread pool divides the $S$ particles across \texttt{num\_threads} workers.
Concretely, the C++ function \texttt{augment\_trees()} accepts an integer
\texttt{num\_threads} argument; setting \texttt{num\_threads = 1} (default)
runs everything sequentially, while any value $>1$ activates the parallel
worker pool. The number of threads should not exceed the number of physical
cores, and diminishing returns are expected when the per-particle compute
time is small relative to thread-scheduling overhead.

% ══════════════════════════════════════════════════════════════════════════════
\section{The Phylogenetic Likelihood}\label{sec:likelihood}

\subsection{Observed versus augmented data}

The observed data $\obs$ consists of the extant tree: branching times
$\brts$ and tip count $n$. The unobserved complete data is the augmented
tree $\aug = (\obs, z)$, where $z$ denotes all extinct lineages. The
marginal likelihood integrates over all possible extinct histories:
\begin{equation}\label{eq:marginal}
  p(\obs \mid \pars) = \int p(\obs, z \mid \pars) \, \dd\nu(z),
\end{equation}
where $\nu$ is the natural dominating measure on augmented-tree space
(counting measure over the combinatorial structure times Lebesgue measure
over the continuous event times). This integral is analytically tractable
for the CR and DD models \citep{etienne2012prolonged} but is intractable
for EP and any model with lineage-specific covariates.

\subsection{Complete-data log-likelihood}\label{sec:loglik}

Let the augmented tree $\aug$ have $K$ total lineages, lineage $k$ alive on
the time interval $[a_k, d_k]$ where $a_k$ is its birth time and $d_k$ its
death (or the crown age for extant lineages). The complete-data
log-likelihood of the birth--death process is:
\begin{equation}\label{eq:loglik}
  \log p(\obs, z \mid \pars)
  = \sum_{k=1}^{K} \log r_k(t_k^{\rm ev})
    - \sum_{k=1}^{K} \int_{a_k}^{d_k}
      \bigl[\lambda(k,t;\pars) + \mu(k,t;\pars)\bigr]\dd t,
\end{equation}
where $r_k(t_k^{\rm ev})$ is the rate of the event terminating lineage $k$
($\lambda$ for speciation, $\mu$ for extinction), and the integral is the
cumulative hazard of lineage $k$ over its lifetime. Expanding:
\begin{align}
  \log p(\obs, z \mid \pars)
  &= \sum_{k:\,\text{speciation}} \log \lambda(k, t_k^{\rm ev};\pars)
   + \sum_{k:\,\text{extinction}} \log \mu(k, t_k^{\rm ev};\pars)
   \nonumber\\
  &\quad - \sum_{k=1}^{K}
    \int_{a_k}^{d_k}
    \bigl[\lambda(k,t;\pars) + \mu(k,t;\pars)\bigr]\dd t.
    \label{eq:loglik2}
\end{align}

\paragraph{Closed-form hazard integrals.}
For the linear link, $\lambda(k,t;\pars) = \max(0, \eta_k(t))$ and the
integral $\int_{a_k}^{d_k} \eta_k(t)\dd t$ has a closed form for all
supported covariate combinations, because $N(t)$, $P(t)$, and $E(k,t)$
are all piecewise-linear between events. The C++ kernel precomputes the
integral for each lineage interval analytically, making evaluation of
\eqref{eq:loglik2} exact and fast.

% ══════════════════════════════════════════════════════════════════════════════
\section{Importance-Sampling Likelihood Estimation}\label{sec:IS}

\subsection{The IS estimator}

Let $z_1, \ldots, z_S \stackrel{\rm iid}{\sim} q(z \mid \obs, \pars)$.
The marginal likelihood \eqref{eq:marginal} is the expectation
\begin{equation}
  p(\obs \mid \pars)
  = \E_{q}\!\left[\frac{p(\obs, z \mid \pars)}{q(z \mid \obs, \pars)}\right],
\end{equation}
estimated without bias by
\begin{equation}\label{eq:IS}
  \Lhat_S(\pars) = \frac{1}{S}\sum_{i=1}^{S}
    \exp\!\bigl(\logf_i - \logq_i\bigr)
    = \frac{1}{S}\sum_{i=1}^{S} \exp(\lw_i).
\end{equation}
The IS log-likelihood estimate is therefore
\begin{equation}\label{eq:fhat}
  \fhat(\pars) = \log \Lhat_S(\pars)
  = \log\!\left(\frac{1}{S}\sum_{i=1}^{S} \exp(\lw_i)\right),
\end{equation}
evaluated in numerically stable form via the log-sum-exp identity:
\begin{equation}\label{eq:logsumexp}
  \fhat(\pars) = w_{\max} + \log\!\left(\frac{1}{S}\sum_{i=1}^{S}
    \exp(\lw_i - w_{\max})\right), \quad
  w_{\max} = \max_i \lw_i.
\end{equation}

\begin{remark}[Jensen's inequality bias]\label{rem:jensen}
Because $\log$ is concave, $\fhat(\pars) \leq \E_q[\lw]$ by Jensen's
inequality, so $\fhat$ underestimates $\log p(\obs\mid\pars)$. The bias
is $O(1/S)$ and vanishes as $S\to\infty$.
\end{remark}

\subsection{Bias-corrected estimator}\label{sec:biascorrect}

A moment-based correction removes the leading-order Jensen bias. Define
$\ell_{\rm hat} = S^{-1}\sum_i \lw_i$ and the central moments
$m_k = S^{-1}\sum_i (\lw_i - \ell_{\rm hat})^k$. The exact identity
relating $\log\E[e^W]$ to the cumulants of $W$ gives:
\begin{equation}\label{eq:biascorrect}
  \fhat_{\rm BC}(\pars) = \ell_{\rm hat}
    + \log\!\left(1 + \sum_{k=2}^{K}\frac{m_k}{k!}\right).
\end{equation}
At $K=2$ this is the log-normal correction $\ell_{\rm hat} + \log(1 + m_2/2)
\approx \ell_{\rm hat} + m_2/2$ for small $m_2$. The full series
($K\to\infty$) recovers $\fhat$ exactly; the truncation at finite $K$ is the
approximation.

\paragraph{Handling rejected trees.}
Trees that exceed $\lambda_{\max}$ or the maximum extinct-lineage count
have $L_i = 0$ ($\lw_i = -\infty$). They are excluded from moment sums
but enter the denominator:
\begin{equation}
  \fhat(\pars) = \fhat_{\rm valid}(\pars)
    + \log\!\left(\frac{S_{\rm valid}}{S}\right),
\end{equation}
where $\fhat_{\rm valid}$ uses only the $S_{\rm valid}$ non-rejected trees.

\subsection{Effective sample size}

The IS effective sample size quantifies IS quality:
\begin{equation}\label{eq:ESS}
  \ESS = \frac{\left(\displaystyle\sum_{i=1}^{S} e^{\lw_i}\right)^2}
              {\displaystyle\sum_{i=1}^{S} e^{2\lw_i}}.
\end{equation}
$\ESS \approx S$ means roughly uniform weights (efficient IS); $\ESS\approx 1$
means one tree dominates (IS collapse). We report $\ESS/S\in(0,1]$ as a
diagnostic. For reliable inference, $\ESS/S \geq 0.1$ is recommended.

\subsection{Cross-importance sampling (Mode~2)}\label{sec:crossIS}

In Mode~2 (\texttt{shared\_trees = TRUE}), a single pool of $S$ augmented
trees is drawn at a reference parameter $\pars_0$, and $\fhat$ is evaluated
at any $\pars_j$ by re-scoring the same trees:
\begin{equation}\label{eq:crossIS}
  \fhat(\pars_j) = \log\!\left(\frac{1}{S}\sum_{i=1}^{S}
    \exp\bigl(\logf(z_i,\pars_j) - \logq_i\bigr)\right),
\end{equation}
where $\logq_i = \log q(z_i \mid \obs, \pars_0)$ are fixed at simulation
time and $\logf(z_i,\pars_j)$ is re-evaluated analytically at $\pars_j$
via \texttt{eval\_logf()}. This reduces the total augmentations from
$M \times S$ to $S$ per CEM iteration, at the cost of higher variance when
$\pars_j$ and $\pars_0$ are distant. Mode~2 is parallelised over the $M$
particles (scoring is cheap), while Mode~1 parallelises the $S$
augmentation draws.

% ══════════════════════════════════════════════════════════════════════════════
\section{The Cross-Entropy Method with Simulated Annealing}\label{sec:CEM}

\subsection{The optimisation problem}

Parameter estimation reduces to maximising the IS log-likelihood:
\begin{equation}\label{eq:opt}
  \hat{\pars} = \argmax_{\pars \in [\lb,\,\ub]} \fhat(\pars).
\end{equation}
Because $\fhat$ is stochastic (it depends on randomly drawn augmented
trees) and potentially multimodal, gradient-based methods are unreliable.
The Cross-Entropy Method (CEM) \citep{rubinstein1999cross,deBoer2005tutorial}
solves \eqref{eq:opt} by maintaining an evolving population of candidate
parameter vectors and adaptively concentrating it near the optimum.

\subsection{Cross-entropy population update}

Let $g(\pars; \bm{\mu}, \bm{\Sigma})$ be a parametric family of densities
over the search box $[\lb, \ub]$. At each iteration the CEM selects the
member of $g$ that minimises the KL divergence to the ``ideal'' distribution
$g^*$ that concentrates all mass at the global maximiser of $\fhat$:
\begin{equation}
  (\bm{\mu}^{(k+1)}, \bm{\Sigma}^{(k+1)}) = \argmin_{\bm{\mu},\bm{\Sigma}}
  \KL\!\left(g^* \;\big\|\; g(\cdot;\bm{\mu},\bm{\Sigma})\right).
\end{equation}
In the Gaussian family this CE update reduces to computing the sample mean
and covariance of the elite particles. Because the full covariance update
can overfit when the elite set is small, \textsf{emphasis} uses a
diagonal approximation: each parameter dimension is updated independently,
giving a product of independent Gaussian perturbations.

Each CEM iteration consists of three steps --- \textbf{Evaluation},
\textbf{Selection}, and \textbf{Update} --- described in the following
subsections and collected in Algorithm~\ref{alg:CEM}. Note that these are
\emph{not} the E-step and M-step of the EM algorithm; the CEM has no
latent-variable expectation and no $Q$-function maximisation.

\subsection{Simulated annealing schedule}\label{sec:annealing}

To prevent premature convergence, the perturbation variance is not derived
purely from the elite sample but is controlled by an explicit
\emph{annealing schedule} analogous to simulated annealing
\citep{kirkpatrick1983optimization}. Concretely, define the ``temperature''
vector $\bm{\sigma}^{(k)} \in \R^8_{>0}$ as the per-parameter perturbation
standard deviation at iteration $k$. New particles are generated as:
\begin{equation}\label{eq:perturb}
  \pars_j^{(k+1)} = \mathrm{clip}\!\left(
    \pars_{\pi(j)}^{(k)} + \bm{\varepsilon}_j,\; \lb,\; \ub\right),
  \quad
  \bm{\varepsilon}_j \sim \mathcal{N}\!\left(\mathbf{0},\;
    \mathrm{diag}\!\left(\bm{\sigma}^{(k)}\right)^2\right),
\end{equation}
where $\pi(j)$ selects uniformly from the elite set and
$\mathrm{clip}(\cdot,\lb,\ub)$ projects back into the feasible box.

The temperature is cooled linearly:
\begin{equation}\label{eq:anneal}
  \bm{\sigma}^{(k+1)} = \bm{\sigma}^{(k)} - \alpha\,\mathbf{1},
  \qquad
  \alpha = \frac{\bm{\sigma}^{(0)}}{T},
\end{equation}
with the initial temperature set to one quarter of the search range:
\begin{equation}\label{eq:sigma0}
  \sigma_p^{(0)} = \frac{u_p - \ell_p}{4}, \quad p = 1,\ldots, 8.
\end{equation}
With this schedule, $\bm{\sigma}^{(T)} = \mathbf{0}$: the population fully
collapses at iteration $T$. The linear decay is analogous to the linear
cooling schedule in classical simulated annealing and guarantees a smooth
transition from broad exploration (large $\sigma$, wide Gaussian
perturbations) to fine-grained exploitation (small $\sigma$, narrow
perturbations near the current elite). Early stopping via the plateau and
annealing rules (Section~\ref{sec:stopping}) allows termination before
$T$ is reached.

\paragraph{Connection to simulated annealing.}
In classical simulated annealing, the Boltzmann acceptance probability at
temperature $\tau$ is $\exp(-\Delta E / \tau)$: high $\tau$ accepts bad
moves freely (exploration); low $\tau$ accepts only improvements
(exploitation). In the CEM-SA framework, $\bm{\sigma}^{(k)}$ plays the role
of $\tau$: large $\sigma$ generates particles far from the current elite
(exploration); as $\sigma \to 0$ only particles close to elite copies are
generated (exploitation). The crucial difference is that CEM-SA evaluates
the stochastic objective $\fhat$ directly at many candidate points in
parallel rather than performing a single-chain random walk.

\subsection{Stopping rules}\label{sec:stopping}

Three criteria are checked in order at the end of each iteration:
\begin{enumerate}
  \item \textbf{Annealing exhausted.} $\bm{\sigma}^{(k)} \leq \mathbf{0}$
    element-wise: the population has fully collapsed; further perturbation
    would be zero.
  \item \textbf{Plateau.} The improvement
    $f^{(k)} - f^{(k-1)} < \delta$ (default $\delta = 10^{-4}$) for $P$
    (default 5) consecutive iterations. This guards against wasting
    computation when the IS surface is flat near the optimum.
  \item \textbf{Hard cap.} $k = T$ (default $T = 20$).
\end{enumerate}

\subsection{Algorithm}\label{sec:CEMalg}

The complete CEM-SA procedure is given in Algorithm~\ref{alg:CEM}.

\begin{algorithm}[H]
\caption{Monte Carlo CEM with Simulated Annealing (\texttt{emphasis\_cem})}
\label{alg:CEM}
\DontPrintSemicolon
\SetKwComment{Comment}{// }{}
\KwIn{branching times $\brts$; bounds $[\lb,\ub]$; population size $M$;
      elite fraction $\rho$; IS sample size $S$; max iterations $T$;
      annealing step $\alpha$; plateau tolerance $\delta$; patience $P$;
      \texttt{num\_threads}}
\KwOut{best parameter estimate $\hat{\pars}$; per-iteration history}
\BlankLine
$\bm{\sigma}^{(0)} \leftarrow (u_p - \ell_p)/4$ for each $p$
  \Comment*[r]{initial temperature \eqref{eq:sigma0}}
Initialise $\{\pars_j^{(0)}\}_{j=1}^{M}$ uniformly in $[\lb,\ub]$\;
$k \leftarrow 0$;\quad \textit{plat} $\leftarrow 0$\;
\BlankLine
\While{$k < T$}{
  \lIf{$\bm{\sigma}^{(k)} \leq \mathbf{0}$}{\textbf{stop}
    (\textit{annealing exhausted})}
  \BlankLine
  \Comment{Step 1 -- Evaluation: score every particle (in parallel)}
  \For(\Comment*[f]{\texttt{num\_threads} workers}){$j = 1, \ldots, M$}{
    Draw $z_1^{(j)},\ldots,z_S^{(j)} \sim q(\cdot \mid \obs, \pars_j^{(k)})$\;
    Compute $\fhat(\pars_j^{(k)})$ via \eqref{eq:fhat};
    mark invalid if result is NA\;
  }
  \BlankLine
  \Comment{Step 2 -- Selection: identify elite particles}
  $j^* \leftarrow \argmax_{j:\,\text{valid}} \fhat(\pars_j^{(k)})$;\quad
  $f^{(k)} \leftarrow \fhat(\pars_{j^*}^{(k)})$\;
  Elite $\leftarrow$ top $\lfloor\rho M\rfloor$ valid particles ranked by $\fhat$\;
  \BlankLine
  \Comment{Plateau stopping check}
  \eIf{$k > 0$ \textbf{and} $f^{(k)} - f^{(k-1)} < \delta$}{
    \textit{plat} $\leftarrow$ \textit{plat} $+ 1$\;
    \lIf{\textit{plat} $\geq P$}{\textbf{stop} (\textit{plateau})}
  }{\textit{plat} $\leftarrow 0$}
  \BlankLine
  \Comment{Step 3 -- Update: resample around elite particles and cool}
  For each non-elite slot $j$: draw $\pi(j) \sim \mathrm{Uniform}(\text{Elite})$
  and generate $\pars_j^{(k+1)}$ via \eqref{eq:perturb}\;
  Copy elite particles unchanged:
  $\pars_j^{(k+1)} \leftarrow \pars_j^{(k)}$ for $j\in\text{Elite}$\;
  $\bm{\sigma}^{(k+1)} \leftarrow \bm{\sigma}^{(k)} - \alpha\,\mathbf{1}$\;
  $k \leftarrow k + 1$\;
}
$\hat{\pars} \leftarrow \argmax_{k} f^{(k)}$
  \Comment*[r]{global best across all iterations}
\end{algorithm}

\paragraph{Parallelism in the evaluation step.}
The inner loop over particles $j = 1,\ldots,M$ is the dominant cost.
Each augmentation is entirely independent: particle $j$ accesses only its
own parameter vector $\pars_j^{(k)}$ and writes to its own pre-allocated
result slot. This embarrassing parallelism is exploited by
\textsf{RcppParallel}: the $M$ particles are partitioned into
\texttt{num\_threads} contiguous blocks, each processed by a separate
thread from the TBB worker pool. Because there is no inter-particle
communication and each thread writes to a disjoint memory region, no locks
or atomic operations are needed, and parallel efficiency is near-linear up
to the number of physical cores.

\paragraph{Particle rescue.}
If every particle in a given iteration produces an invalid $\fhat$ (all
augmented trees rejected), the per-particle limits are escalated:
$\lambda_{\max} \leftarrow 10\lambda_{\max}$,
$N_{\max} \leftarrow 10 N_{\max}$,
and the E-step is repeated. If $N_{\max} > 10{,}000$ with all particles
still invalid, the run terminates with a warning advising the user to
increase \texttt{num\_points} or widen the bounds.

% ══════════════════════════════════════════════════════════════════════════════
\section{Monte Carlo EM for Local Refinement}\label{sec:MCEM}

\subsection{The EM objective}

The Monte Carlo Expectation--Maximisation (MCEM) algorithm
\citep{wei1990monte,levine2001implementations} maximises the expected
complete-data log-likelihood:
\begin{equation}\label{eq:Qfunc}
  Q(\pars \mid \pars^{(t)}) = \E_{q(\cdot\mid\obs,\pars^{(t)})}\!
    \bigl[\log p(\obs, z \mid \pars)\bigr].
\end{equation}

\subsection{IS-weighted E-step}

At iteration $t$, draw $z_1,\ldots,z_S \sim q(\cdot\mid\obs,\pars^{(t)})$.
The IS-weighted approximation to $Q$ is:
\begin{equation}\label{eq:Qhat}
  \hat{Q}(\pars \mid \pars^{(t)})
  = \frac{\displaystyle\sum_{i=1}^{S}
          \tilde{w}_i^{(t)} \cdot \logf(z_i, \pars)}
         {1},
  \quad
  \tilde{w}_i^{(t)} = \frac{e^{\lw_i^{(t)}}}{\sum_j e^{\lw_j^{(t)}}},
\end{equation}
where $\lw_i^{(t)} = \logf(z_i,\pars^{(t)}) - \logq_i^{(t)}$ and
$\tilde{w}_i^{(t)}$ are normalised IS weights. This is the self-normalised
IS estimator of $Q$, which is consistent for any positive proposal
$q > 0$ wherever $p > 0$.

\subsection{M-step and convergence}

The M-step maximises $\hat{Q}(\pars\mid\pars^{(t)})$ over $\pars\in[\lb,\ub]$
using the derivative-free COBYLA algorithm \citep{powell1994direct} via the
\texttt{nloptr} interface. COBYLA constructs successive linear approximations
to the objective from function evaluations, which is appropriate here because
the IS noise makes numerical gradients unreliable.

Convergence is declared when the change in $\fhat(\pars^{(t)})$ across
consecutive post-burnin iterations falls below a tolerance $\varepsilon$
(default 0.1). A burn-in phase (default 20 iterations) is excluded from
the convergence check to allow transients to decay.

\subsection{MCEM as a refinement stage}

Because MCEM converges to a local maximum of $\pars \mapsto
\log p(\obs\mid\pars)$, the recommended workflow is:
\begin{enumerate}
  \item Run CEM-SA to locate a good region (global search).
  \item Warm-start MCEM at $\hat{\pars}_{\rm CEM}$ (local refinement).
\end{enumerate}
The MCEM E-step parallelises over the $S$ augmented trees in the same way
as the CEM E-step, using \texttt{num\_threads} threads.

% ══════════════════════════════════════════════════════════════════════════════
\section{Model Selection}\label{sec:modelsel}

\subsection{AIC}

For fitted model $\mathcal{M}_k$ with $p_k$ free parameters and IS
log-likelihood $\fhat(\hat{\pars}_k)$:
\begin{equation}\label{eq:AIC}
  \AIC_k = -2\,\fhat(\hat{\pars}_k) + 2 p_k.
\end{equation}
The preferred model minimises $\AIC$. The gaps
$\Delta\AIC_k = \AIC_k - \min_j \AIC_j \geq 0$ quantify evidence against
model $k$: $\Delta\AIC < 2$ is considered negligible,
$2 \leq \Delta\AIC < 6$ moderate, and $\Delta\AIC \geq 6$ strong.

\begin{remark}[IS bias in AIC comparisons]\label{rem:aicbias}
Because $\fhat$ underestimates $\log p(\obs\mid\hat{\pars})$
(Remark~\ref{rem:jensen}), the AIC is upward-biased. When models are
evaluated at independently drawn augmented trees (Mode~1), the bias is
model-specific and does not cancel in $\Delta\AIC$. This residual Monte
Carlo noise is quantified by the bootstrap variance
(Section~\ref{sec:bootvar}) and is reduced by increasing $S$ or using
Mode~2.
\end{remark}

\subsection{Bootstrap variance of $\fhat$}\label{sec:bootvar}

To quantify IS Monte Carlo uncertainty at the optimum, \textsf{emphasis}
provides a leave-half-out subsampling bootstrap:
\begin{equation}\label{eq:bootvar}
  \widehat{\Var}[\fhat(\hat{\pars})] = \frac{1}{B} \sum_{b=1}^{B}
    \bigl(\fhat^{(b)}(\hat{\pars}) - \bar{\fhat}(\hat{\pars})\bigr)^2,
\end{equation}
where $\fhat^{(b)}$ is computed on a uniformly drawn half of the
$S_{\rm final}$ trees and $\bar{\fhat}$ is the mean over $B$ replicates.
If $\widehat{\Var}^{1/2}[\fhat] \sim |\Delta\AIC|/2$, the model comparison
is noise-dominated and $S$ should be increased.

\subsection{Automated three-model selection}

\texttt{select\_diversification\_model()} fits the three biologically most
relevant models (DD, PD, EP) in a single call, applies \eqref{eq:AIC}, and
returns a ranked comparison table with AIC weights
$w_k = e^{-\Delta\AIC_k/2} / \sum_j e^{-\Delta\AIC_j/2}$.

% ══════════════════════════════════════════════════════════════════════════════
\section{Diagnostics}\label{sec:diagnostics}

\subsection{CEM convergence diagnostics}

\texttt{diagnose\_cem()} extracts and plots four diagnostic panels:
\begin{description}
  \item[Log-likelihood trace.] $f^{(k)}$ vs.\ iteration $k$, annotated
    with the stopping rule. A monotone increasing curve ending at a plateau
    indicates healthy convergence.
  \item[Parameter traces.] $\hat{\pars}^{(k)}_p$ vs.\ $k$ for each
    dimension $p$. Stable late-iteration values confirm convergence.
  \item[Population spread.] Median and IQR of $\fhat(\pars_j^{(k)})$ over
    all valid particles. Narrowing IQR reflects progressive collapse.
  \item[IS weight distribution.] Histogram of $\lw_i$ at the globally
    best particle, annotated with $\ESS$. A concentrated unimodal
    histogram near the mean indicates healthy IS weights.
\end{description}

\subsection{GAM-based survival probability}

\texttt{train\_GAM()} fits a binomial GAM to batch-simulation outcomes,
estimating $\pi(\pars) = \Pr(\text{clade survives}\mid\pars)$ as a smooth
surface. \texttt{predict\_survival()} evaluates this surface at any $\pars$
without re-simulating, enabling fast likelihood corrections for conditioning
on non-extinction.

% ══════════════════════════════════════════════════════════════════════════════
\section{Implementation}\label{sec:impl}

\subsection{Software architecture}

\textsf{emphasis} is written in \textsf{R} and \textsf{C++17}, with the
computationally intensive steps delegated via \textsf{Rcpp}. The public API
consists of five primary functions:
\begin{description}
  \item[\texttt{simulate\_tree()}] Forward and conditional simulation;
    dispatches to the C++ Gillespie/thinning engine.
  \item[\texttt{augment\_trees()}] E-step primitive: draw $S$ augmented trees,
    return \texttt{trees}, \texttt{logf}, \texttt{logg}.
  \item[\texttt{eval\_logf()}] Re-score cached trees at a new $\pars$
    (enables Mode~2 cross-IS, Section~\ref{sec:crossIS}).
  \item[\texttt{estimate\_rates()}] Unified inference entry point; dispatches
    to CEM-SA or MCEM and returns a named \texttt{emphasis\_fit} object.
  \item[\texttt{compare\_models()} / \texttt{select\_diversification\_model()}]\mbox{}\\
    AIC-based model comparison.
\end{description}

\subsection{Separation of simulation and scoring}

The key architectural principle is the clean separation of three concerns:
\begin{enumerate}
  \item \textbf{Simulation.} \texttt{augment\_trees()} simulates trees and
    computes $\logf$ and $\logq$ in a single pass.
  \item \textbf{Scoring.} \texttt{eval\_logf()} re-evaluates $\logf(z,\pars')$
    for a new $\pars'$ on already-simulated trees, without re-running the
    stochastic augmentation.
  \item \textbf{IS aggregation.} \texttt{.is\_fhat(logf, logg)} in \textsf{R}
    combines $\logf$ and $\logq$ into $\fhat$ via \eqref{eq:fhat}.
\end{enumerate}
This design is what makes Mode~2 possible: one simulation pass at $\pars_0$
produces trees reused across all $M$ particles.

\subsection{Parallel computing}\label{sec:parallel}

All parallelism in \textsf{emphasis} is shared-memory multithreading
implemented via \textsf{RcppParallel}, which provides:
\begin{itemize}
  \item A \emph{portable thread pool} backed by Intel TBB
    (\texttt{tbb::parallel\_for}) on Linux, macOS, and Windows, with a
    fallback \texttt{std::thread} implementation when TBB is unavailable.
  \item \emph{Thread-safe memory} via pre-allocated per-thread result
    vectors: each thread writes exclusively to its own slice of the output
    array with no shared mutable state, so no locks or atomics are needed.
  \item \emph{Work scheduling}: the $S$ (Mode~1) or $M$ (Mode~2) tasks
    are divided into \texttt{num\_threads} contiguous blocks. The block
    boundaries are fixed at launch time, so there is no dynamic work
    stealing (which keeps overhead minimal for uniform per-task cost).
\end{itemize}

\paragraph{Practical guidance.}
\begin{itemize}
  \item Set \texttt{num\_threads} to the number of \emph{physical} cores
    (not hyperthreads); hyperthreading rarely helps for floating-point
    intensive workloads.
  \item Parallel speedup is near-linear when $S \gg \texttt{num\_threads}$
    and per-tree compute time dominates thread-launch overhead. For very
    small $S$ (e.g.\ $S=1$, the CEM default), single-threaded mode is
    optimal.
  \item In Mode~1, parallelise the $S$ trees per particle
    (\texttt{num\_threads} $> 1$ inside \texttt{augment\_trees}).
    In Mode~2, the $M$ re-scorings in \texttt{eval\_logf} are parallelised
    instead, since augmentation itself is already done.
\end{itemize}

\subsection{Parameter layout}

All C++ functions operate on the full 8-element parameter vector in the
canonical order
$(\beta_0, \beta_N, \beta_P, \beta_E, \gamma_0, \gamma_N, \gamma_P,
\gamma_E)^\top$.
Inactive parameters are set to zero. The model binary flag
$\mathbf{m} = (m_N, m_P, m_E) \in \{0,1\}^3$ gates which covariate terms
are included in the augmentation proposal. For the EP covariate, $m_E = 1$
additionally switches the C++ rate functions from standard to EP-aware
mode in \texttt{eval\_logf}, since the EP contribution to the hazard
integral \eqref{eq:hazard} requires a different analytic formula.

% ══════════════════════════════════════════════════════════════════════════════
\section{Simulation Studies}\label{sec:simstudies}

\subsection{Estimation power}\label{sec:study1}

To evaluate the frequentist accuracy of CEM estimation under the PD model,
we simulate $n = 100$ independent trees from the true parameter vector
\[
  \pars^* = (\beta_0, \beta_P, \gamma_0, \gamma_P)
          = (0.6,\; 0.005,\; 0.15,\; -0.003),
\]
with crown age $T = 8$ and estimate $\pars$ on each replicate via CEM with
$M = 100$, $S = 1$, $T = 20$ iterations. For each replicate $r$ we record
the estimate $\hat{\pars}_r$ and the number of extant tips $n_r$.
Accuracy is summarised by:
\begin{align}
  \text{Bias}_p &= \frac{1}{n}\sum_{r=1}^{n}
    (\hat{\pars}_{r,p} - \pars^*_p),\label{eq:bias}\\
  \text{RMSE}_p &= \sqrt{\frac{1}{n}\sum_{r=1}^{n}
    (\hat{\pars}_{r,p} - \pars^*_p)^2}.\label{eq:rmse}
\end{align}
The tree-size effect is visualised by plotting $|\hat{\pars}_{r,p} - \pars^*_p|$
against $n_r$ with a LOWESS smooth. Monotone decrease in error with $n_r$
confirms consistency: larger trees produce a sharper IS likelihood surface.

\subsection{Model selection power}\label{sec:study2}

We assess AIC-based model selection in a fully crossed design. Twenty-five
trees are simulated from each of the four pure models (CR, DD, PD, EP),
giving 100 trees total. All four models are fitted to each tree via CEM,
and the model with the lowest AIC is selected. Results are summarised by
the $4\times 4$ confusion matrix
\begin{equation}
  C_{ij} = \#\bigl\{\text{trees from model }i
             \text{ for which model }j\text{ is selected}\bigr\},
\end{equation}
overall accuracy $\mathrm{acc} = \operatorname{tr}(C)/\sum_{i,j}C_{ij}$,
and per-model recall $C_{ii}/\sum_j C_{ij}$.

% ══════════════════════════════════════════════════════════════════════════════
\section{Summary of Control Parameters}\label{sec:control}

\begin{table}[ht]
\centering
\caption{CEM control parameters (\texttt{estimate\_rates\_control("cem")}).}
\label{tab:cemcontrol}
\renewcommand{\arraystretch}{1.2}
\begin{tabular}{llp{7.8cm}}
\toprule
Parameter & Default & Description \\
\midrule
\texttt{max\_iter} & 20 & Hard iteration cap $T$ \\
\texttt{num\_points} & 50 & Population size $M$ \\
\texttt{sample\_size} & 1 & IS sample size $S$ per particle per iteration \\
\texttt{shared\_trees} & \texttt{FALSE} & Pool trees across particles (Mode~2) \\
\texttt{disc\_prop} & 0.5 & Elite fraction $\rho$ \\
\texttt{sd\_vec} & \texttt{NULL} & Initial $\bm{\sigma}^{(0)}$; auto $=(u-\ell)/4$ \\
\texttt{maxN} & 10 & Max augmentation attempts before rejection \\
\texttt{max\_missing} & $10^4$ & Max extinct lineages per augmentation \\
\texttt{max\_lambda} & 500 & Thinning bound $\lambda_{\max}$ \\
\texttt{tol} & $10^{-4}$ & Plateau tolerance $\delta$ \\
\texttt{patience} & 5 & Consecutive plateau iterations $P$ \\
\texttt{bias\_correct} & \texttt{FALSE} & Moment-based bias correction \eqref{eq:biascorrect} \\
\texttt{n\_boot} & 0 & Bootstrap replicates for $\widehat{\Var}[\fhat]$ \\
\texttt{num\_threads} & 1 & Parallel threads (Section~\ref{sec:parallel}) \\
\bottomrule
\end{tabular}
\end{table}

\begin{table}[ht]
\centering
\caption{MCEM control parameters (\texttt{estimate\_rates\_control("mcem")}).}
\label{tab:mcemcontrol}
\renewcommand{\arraystretch}{1.2}
\begin{tabular}{llp{7.8cm}}
\toprule
Parameter & Default & Description \\
\midrule
\texttt{sample\_size} & 200 & IS sample size $S$ per EM iteration \\
\texttt{tol} & 0.1 & Convergence tolerance $\varepsilon$ \\
\texttt{burnin} & 20 & Iterations excluded from convergence check \\
\texttt{max\_missing} & $10^4$ & Max extinct lineages per augmentation \\
\texttt{max\_lambda} & 500 & Thinning bound $\lambda_{\max}$ \\
\texttt{num\_threads} & 1 & Parallel threads (Section~\ref{sec:parallel}) \\
\bottomrule
\end{tabular}
\end{table}

% ══════════════════════════════════════════════════════════════════════════════
\section{Mathematical Appendix}\label{sec:appendix}

\subsection{Log-sum-exp identity}

For numerical stability \eqref{eq:fhat} is evaluated as:
\begin{align}
  \log\!\left(\frac{1}{S}\sum_{i=1}^S e^{\lw_i}\right)
  &= \log\!\left(\frac{e^{w_{\max}}}{S}
      \sum_{i=1}^S e^{\lw_i - w_{\max}}\right)
   = w_{\max} + \log\!\left(\frac{1}{S}\sum_{i=1}^S
     e^{\lw_i - w_{\max}}\right).
\end{align}
Since $\lw_i - w_{\max} \leq 0$, all terms $e^{\lw_i-w_{\max}} \in (0,1]$
are safe.

\subsection{ESS in log-space}

Define normalised weights $\tilde{w}_i = e^{\lw_i}/\sum_j e^{\lw_j}$.
Then $\ESS = 1/\sum_i \tilde{w}_i^2$, which in log-space is:
\begin{equation}
  \ESS = \exp\!\bigl[2\,\mathrm{lse}(\mathbf{w})
                   - \mathrm{lse}(2\mathbf{w})\bigr],
\end{equation}
where $\mathrm{lse}(\mathbf{v}) = \max_i v_i
+ \log\sum_i e^{v_i - \max_j v_j}$.

\subsection{Annealing decay rate}

With $\bm{\sigma}^{(0)} = \bm{\sigma}_0$ and linear decay
$\bm{\sigma}^{(k)} = \bm{\sigma}_0 - k\alpha\mathbf{1}$, the population
collapses ($\sigma_p^{(k^*)} \leq 0$) at $k^* = \lceil\sigma_{0,p}/\alpha\rceil$.
Setting $k^* = T$ for all $p$ gives $\alpha = \sigma_{0,p}/T$ as in
\eqref{eq:anneal}. This ensures the annealing-exhausted stopping rule
\emph{would} trigger exactly at iteration $T$ in the absence of plateau
stopping.

\subsection{Hazard integral for the linear link}\label{sec:hazard}

For lineage $k$ alive on $[a_k, d_k]$ under the linear link, the hazard
integral decomposes as:
\begin{equation}
  H_k = \int_{a_k}^{d_k}
    \bigl[f(\eta_\lambda(k,t)) + f(\eta_\mu(k,t))\bigr]\dd t,
\end{equation}
where $\eta_\lambda(k,t) = \beta_0 + \beta_N N(t) + \beta_P P(t)
+ \beta_E E(k,t)$ and similarly for $\eta_\mu$. Since $N(t)$, $P(t)$,
and $E(k,t)$ are all piecewise-linear in $t$ between consecutive events,
each sub-interval $[t_j, t_{j+1}]$ within $[a_k, d_k]$ has a constant
$N$, $P$, and linearly increasing $E(k,t) = t - a_k$. On such a
sub-interval with $f(\eta) = \max(0,\eta)$:
\begin{equation}
  \int_{t_j}^{t_{j+1}} \max(0,\, c_0 + c_1 t)\dd t
  = \begin{cases}
      c_0(t_{j+1}-t_j) + \tfrac{c_1}{2}(t_{j+1}^2 - t_j^2)
        & \text{if } c_0 + c_1 t \geq 0 \text{ on } [t_j,t_{j+1}], \\
      0 & \text{if } c_0 + c_1 t \leq 0 \text{ on } [t_j,t_{j+1}], \\
      \text{partial integral} & \text{otherwise (split at root).}
    \end{cases}
\end{equation}
The C++ kernel evaluates this analytically for each sub-interval, making
the complete-data log-likelihood \eqref{eq:loglik2} exact.

% ══════════════════════════════════════════════════════════════════════════════
\bibliographystyle{plainnat}
\begin{thebibliography}{9}

\bibitem[de~Boer et~al.(2005)]{deBoer2005tutorial}
de~Boer, P.-T., Kroese, D.~P., Mannor, S., and Rubinstein, R.~Y. (2005).
A tutorial on the cross-entropy method.
\textit{Annals of Operations Research}, 134(1):19--67.

\bibitem[Etienne et~al.(2012)]{etienne2012prolonged}
Etienne, R.~S., Haegeman, B., Stadler, T., Aze, T., Pearson, P.~N.,
Purvis, A., and Phillimore, A.~B. (2012).
Diversity-dependence brings molecular phylogenies closer to agreement with
the fossil record.
\textit{Proceedings of the Royal Society B}, 279(1732):1300--1309.

\bibitem[Gillespie(1977)]{gillespie1977exact}
Gillespie, D.~T. (1977).
Exact stochastic simulation of coupled chemical reactions.
\textit{The Journal of Physical Chemistry}, 81(25):2340--2361.

\bibitem[Janzen et~al.(2015)]{janzen2015diversity}
Janzen, T., H\"{o}hna, S., and Etienne, R.~S. (2015).
Approximate Bayesian computation of diversification rates from molecular
phylogenies: introducing a new efficient summary statistic, the nLTT.
\textit{Methods in Ecology and Evolution}, 6(5):566--575.

\bibitem[Kirkpatrick et~al.(1983)]{kirkpatrick1983optimization}
Kirkpatrick, S., Gelatt, C.~D., and Vecchi, M.~P. (1983).
Optimization by simulated annealing.
\textit{Science}, 220(4598):671--680.

\bibitem[Levine and Casella(2001)]{levine2001implementations}
Levine, R.~A. and Casella, G. (2001).
Implementations of the Monte Carlo EM algorithm.
\textit{Journal of Computational and Graphical Statistics}, 10(3):422--439.

\bibitem[Lewis and Shedler(1979)]{lewis1979simulation}
Lewis, P.~A.~W. and Shedler, G.~S. (1979).
Simulation of nonhomogeneous Poisson processes by thinning.
\textit{Naval Research Logistics Quarterly}, 26(3):403--413.

\bibitem[Powell(1994)]{powell1994direct}
Powell, M.~J.~D. (1994).
A direct search optimization method that models the objective and constraint
functions by linear interpolation.
In Gomez, S. and Hennart, J.-P. (Eds.),
\textit{Advances in Optimization and Numerical Analysis}, pp.~51--67.
Kluwer Academic Publishers.

\bibitem[Rubinstein(1999)]{rubinstein1999cross}
Rubinstein, R.~Y. (1999).
The cross-entropy method for combinatorial and continuous optimization.
\textit{Methodology and Computing in Applied Probability}, 1(2):127--190.

\bibitem[Wei and Tanner(1990)]{wei1990monte}
Wei, G.~C.~G. and Tanner, M.~A. (1990).
A Monte Carlo implementation of the EM algorithm and the poor man's data
augmentation algorithms.
\textit{Journal of the American Statistical Association}, 85(411):699--704.

\end{thebibliography}

\end{document}
